\section{The full $\operatorname{Sp}_{2n}(2)$ core family}
\label{sec:higher-even-symplectic-draft}

\noindent\textbf{Status.} Author derivation and finite certificates complete. Independent whole-proof
review by /root/g34\_eightfolds passes conditionally precisely on (GM);
the reviewer separately reconstructed the original two-table finite input.
The unconditional three-table endpoint exclusion is stated in D-HES-F below. The proposed all-rank conclusion below
has one explicitly unresolved source-verification gate, denoted (GM): the exact
proper-parabolic semisimple estimate from Gluck--Magaard (1998), Theorem~3.1.
Its statement and finite exception list have been located in the zbMATH review,
and the original bibliographic identity has been verified through Crossref and
the author's institutional repository. Neither that review nor bibliographic
metadata is being promoted to a verified primary mathematical input. The
original article or an exact primary research restatement must be read before
the conclusion is integrated as an unconditional theorem. The elementary
arguments and finite statements below do not depend on guessing that source.

\subsection*{D-HES-01. The finite endpoint and elementary degree bounds}

Put $S_n=\operatorname{Sp}_{2n}(2)$ for $n\ge4$. These groups are simple,
have trivial Schur multiplier and have no outer automorphisms. The complete
exceptional-multiplier list in \cite[Chapter~5, Table~1, numbered p.~35]{BI:HHM}
has only $\operatorname{Sp}_4(2)'$ and $\operatorname{Sp}_6(2)$ in this
even symplectic family. The exact $S_n$, $n\ge4$, multiplier and automorphism
statements are also recorded in Gluck--Magaard, \emph{Absolute Fourth Moments},
author p.~25. These classification statements are used as known inputs; their
truth is independent of that paper's small-fourth-moment hypothesis.

\begin{lemma}[D-HES-01: exact endpoint]
Every nontrivial irreducible complex character of $S_4$ has degree at least
$35$ and satisfies
\[
                   |\chi(x)|\le \frac35\chi(1)\quad(x\ne1).
\tag{HES.1}
\]
Every nontrivial irreducible complex character of $S_5$ has degree at least
$155$. For $n\ge5$, every nontrivial irreducible complex representation of
$S_n$ has degree
\[
                  a\ge \max\{155,2^{n-1}-1\}.
\tag{HES.2}
\]
\end{lemma}

\begin{proof}
The finite assertions are certified from the complete ordinary CTblLib
tables \texttt{S8(2)}, \texttt{S10(2)} and \texttt{S12(2)} supplied by CTblLib. The retained
input contains every class and every irreducible character, with sizes, element
orders, power maps and full provenance. The respective table sizes are $81$, $198$ and $477$. The first two have ATLAS provenance; the third is attributed
by CTblLib to Christoph K\"ohler using group-theoretic and
character-theoretic methods, without a separately verified publication
claim here. The exporter checks full row orthogonality and reality in exact
cyclotomic arithmetic. A separate Python checker reads only the frozen input.
For an entry written as $\sum_j c_j\zeta_N^j$ it uses the rigorous estimate
$|\sum_j c_j\zeta_N^j|\le\sum_j|c_j|$, and compares only rational numbers.
It checks all $271785$ nonidentity entries in nontrivial rows. For $S_4$ the
largest bound is exactly $3/5$, attained at the rational value $-21$ in the
degree-$35$ row, class $2a$. As a supplementary check, for $S_5$ the largest
bound is $17/31$, attained at $-85$ in degree $155$; this stronger ratio is
not required for the all-rank argument. The third table gives minimum degree
$651$ and maximum ratio $11/21$, attained at $-341$ in that degree. Thus all
three tables satisfy the unconditional common ratio bound $3/5$. The code also independently checks
the group-order formulas, class-size sums, degree-square sums, all character
degrees, reality of every entry and the symbolic inequalities used below.

The natural block embedding $S_5\hookrightarrow S_n$ gives $a\ge155$:
a nontrivial irreducible representation of the simple centerless group
$S_n$ is faithful, so its restriction to $S_5$ has a nontrivial constituent.
For the other bound, embed $\operatorname{GL}_n(2)$ in $S_n$ by
$A\mapsto\operatorname{diag}(A,A^{-\mathsf T})$. Its subgroup
\[
 U=\left\{\begin{pmatrix}1&v\\0&I_{n-1}\end{pmatrix}:
                  v\in\mathbf F_2^{n-1}\right\}
\]
is elementary abelian. Its normalizer contains $\operatorname{GL}_{n-1}(2)$,
which is transitive on the $2^{n-1}-1$ nontrivial complex characters of $U$.
Faithfulness makes some such weight occur; conjugation makes all of them
occur with the same positive multiplicity. Thus $a\ge2^{n-1}-1$.
\end{proof}

\noindent\textbf{Finite files.}
\texttt{export\_sp2n2.g}, \texttt{sp2n2\_character\_inputs.json},
\texttt{check\_sp2n2.py}, \texttt{check\_sp2n2.json}.
Frozen input SHA-256:
\texttt{5d3539d84b4988ac34c2703cbc81369b2840e50e3d74ea3afbb2db31f974e615}.
The exact Fourier helper imported by the checker is the already reviewed
\texttt{sporadic/verify\_quarter\_bound.py}; here it is used only for
cyclotomic polynomials in the reality check, not for selecting rows.

\subsection*{D-HES-02. A prime-order eigenspace bound}

The following is the precise source gate, restricted to the case needed here:
\[
\begin{gathered}
\text{(GM): If } n\ge5,\ h\in S_n\text{ is nonidentity semisimple and}\
\text{lies in a proper parabolic subgroup, then every nontrivial}\
\chi\in\operatorname{Irr}(S_n)\text{ satisfies }
                         |\chi(h)|\le\tfrac56\chi(1).
\end{gathered}
\tag{HES.GM}
\]
The review locates this in Gluck--Magaard, \emph{Cross-Characteristic
Character and Fixed Point Space Ratios for Groups of Lie Type}, J. Algebra
\textbf{204} (1998), 188--201, Theorem~3.1(b),
doi:10.1006/jabr.1997.7277. Its displayed exceptional families include
$B_3(2),B_4(2)$, but not $B_n(2)$ for $n\ge5$. If the primary statement
is formulated for cross-characteristic Brauer characters, take a prime not
dividing $|S_n|$, so that the ordinary characters are included. This paragraph
does not substitute for verification of the precise primary hypotheses.

Two verified known inputs used separately are Gow's theorem that every
element of an even-characteristic symplectic group is conjugate to its
inverse, and Gluck's unipotent estimate
\[
               |\chi(u)|\le\tfrac34\chi(1)\quad(1\ne u\text{ unipotent}).
\tag{HES.3}
\]
The latter is explicitly stated in Liebeck's published survey, Theorem~2.4(i),
author pp.~3--4, with the original reference Gluck (1995). The nonunipotent
part of that theorem has bound $19/20$, and is \emph{not} being replaced by
$3/4$ here.

\begin{lemma}[D-HES-02: conditional all-element support]
Assume (GM). If $n\ge5$, $\rho$ is a nontrivial irreducible complex
representation of $S_n$ of degree $a$, and $1\ne x\in S_n$, then
\[
 a-\max_\lambda\dim\ker(\rho(x)-\lambda I)\ge a/9.
\tag{HES.4}
\]
For $n=4,5,6$ the unconditional stronger bound is $a/5$, by the full finite tables.
\end{lemma}

\begin{proof}
We first handle an odd-prime-order element $h$ which belongs to no proper
parabolic. Let $W$ be the natural $2n$-dimensional symplectic space over
$\mathbf F_2$. Since $h$ is semisimple, decompose $W$ into its primary
summands. A summand for an irreducible polynomial not equal to its reciprocal
is totally isotropic: the symplectic pairing pairs it only with the distinct
reciprocal summand. Such a summand would put $h$ in its isotropic-space
stabilizer, a proper parabolic. Likewise a nonzero $1$-eigenspace is a
nondegenerate symplectic space and contains an $h$-fixed isotropic line.
Thus neither type occurs.

Every remaining irreducible polynomial is self-reciprocal of degree $2d$.
On an isotypic summand put $E=\mathbf F_{2^{2d}}$ and
$E_0=\mathbf F_{2^d}$. Adjointness for the symplectic pairing acts on $E$
by its nontrivial involution $z\mapsto\bar z=z^{2^d}$. The nondegenerate
trace pairing uniquely writes the original pairing as the trace of a
nondegenerate Hermitian form $H$ on an $E$-space: more explicitly define
$H(v,w)$ by $B(av,w)=\operatorname{Tr}_{E/\mathbf F_2}(aH(v,w))$ for all
$a\in E$. The adjoint identity and symmetry of $B$ give
$H(av,bw)=a\bar bH(v,w)$ and $H(w,v)=\overline{H(v,w)}$.
If its $E$-dimension is at least two, diagonalization and surjectivity of
$E^\times\to E_0^\times$, $z\mapsto z\bar z$, produce a nonzero isotropic
vector. Its entire $E$-line is totally isotropic for $B$, again a forbidden
invariant isotropic space. Hence each primary polynomial occurs once.

Consequently the symplectic centralizer is a product of norm-one scalar
groups, of orders $2^{d_i}+1$, with $\sum_i d_i=n$. In particular
\[
                  |C_{S_n}(h)|=\prod_i(2^{d_i}+1)\le3^n.
\tag{HES.5}
\]
Character column orthogonality gives
$|\chi(h)|\le|C_{S_n}(h)|^{1/2}$. Combining (HES.2) and (HES.5) yields
\[
       |\chi(h)|/a\le 3^{n/2}/\max\{155,2^{n-1}-1\}<5/6.
\tag{HES.6}
\]
For $n=5,6$ this is the immediate comparison with $155$. For $n\ge7$ square
the inequality; the initial case is
$36\cdot3^7<25\cdot63^2$. If $x=2^{n-1}\ge64$, then
$(2x-1)^2-3(x-1)^2=x^2+2x-2>0$, proving the induction. This supplies
the all-rank estimate without extrapolating any finite character table.

If instead $h$ belongs to a proper parabolic, apply (GM). Every nonidentity
power of a prime-order $h$ generates the same cyclic group, so the same
parabolic alternative and the same bounds apply to all those powers.
For odd prime order $p$, the $1$-eigenvalue multiplicity is therefore at most
\[
 \frac1p\left(a+\sum_{j=1}^{p-1}|\chi(h^j)|\right)
 \le a\left(\frac56+\frac1{6p}\right)\le\frac89a.
\tag{HES.7}
\]
Gow's theorem makes $\chi$ real valued. Thus each nonreal eigenvalue has the
same multiplicity as its complex conjugate; every eigenvalue other than $1$
of this odd-order operator has multiplicity at most $a/2$. This proves the
projective support bound $a/9$ for odd-prime-order elements. If $h$ has order
two, it is unipotent, and (HES.3) gives the stronger bound $a/8$ for both
eigenvalues. For arbitrary $1\ne x$, choose a nonidentity prime-order power
$h=x^k$. Each $x$-eigenspace lies in an $h$-eigenspace, so the projective
support of $x$ is at least that of $h$. Finally the finite common character-ratio bound $3/5$ and the triangle
inequality give the unconditional bound $a/5$ for $S_4,S_5,S_6$.
\end{proof}

\subsection*{D-HES-03. Full normalizers, tensors and the actual Hurwitz budget}

\begin{proposition}[D-HES-03: conditional exclusion of the whole family]
Assume (GM). In the signature-$(3,3)$ sixfold-source problem, an irreducible
inducing representation whose image has an irreducible normal subgroup
which is a product of groups $\operatorname{Sp}_{2n_i}(2)$, $n_i\ge4$,
cannot yield a period image of dimension nine. Arbitrary normalizing
scalars, compatible permutations of factors, and inducing character fields
larger than the original imaginary quadratic field are allowed.
\end{proposition}

\begin{proof}
Let $J\le\operatorname{GL}(V)$ be the faithful inducing image, $e=\dim V$,
and let $E$ be the specified irreducible normal subgroup. Since its factors
are centerless, it is an actual direct product. Complete reducibility gives
$V=\bigotimes_i V_i$, with each $V_i$ a nontrivial irreducible module of a
factor. An operator normalizing $E$ permutes isomorphic simple factors.
After removing this permutation, its restrictions to individual factors
are inner automorphisms, since all these outer automorphism groups are
trivial. Schur's lemma then identifies the operator with a scalar times a
tensor of core operators. There is no projective-cover ambiguity: the Schur
multipliers are trivial, and the restrictions are actual representations.

On a tensor of factor-preserving operators, the largest eigenvalue
multiplicity divided by $e$ is no larger than the corresponding largest
ratio on any nonscalar factor (convolution of probability distributions).
For a nontrivial factor-permutation cycle, conjugating the tensor coordinates
puts all twists on one coordinate, then diagonalizes their cyclic product.
Basis words which are not constant have orbit size at least two, while
there are only $a$ constant words among $a^l$ words for an $l$-cycle.
Every eigenvalue multiplicity is therefore at most
$a+(a^l-a)/2\le3a^l/4$, since $a\ge2$ and $l\ge2$.
Other tensor factors cannot increase this normalized maximum. Thus every
nonscalar element of $J$ has projective support at least $e/9$; if all factors
are among $S_4,S_5,S_6$, the stronger bound $e/5$ holds. A nonidentity scalar has actual
moving rank $e$, so it obeys both lower bounds for actual moving rank.

Gow's theorem makes each $V_i$ self-dual, hence $V$ has a nondegenerate
$E$-invariant bilinear form whose line is unique. Every element of $J$
preserves this line, defining a multiplier character
$\mu:J\to\mathbf C^\times$. It is nontrivial in our problem. Indeed, if
$\mu=1$, then every finite-order $J$-operator has eigenvalues in reciprocal
pairs and hence real trace. The character field of the inducing
representation would be contained in $\mathbf R$, contrary to its containing
the original imaginary quadratic field. This uses field inclusion only;
no equality with that field is assumed.

The actual intermediate-cover argument already established in D-BI-01
identifies the six-dimensional complex multiplicity variation with the
parabolic cohomology of the inducing representation. Its Euler formula is
\[
                 6=2(h-1)e+\sum_{j=1}^r\operatorname{rank}(g_j-I).
\tag{HES.8}
\]
Here the image is faithful and $g_j\ne1$ are its effective inertias.
Since $e\ge35$, $h\ge2$ is impossible. For $h=1$, (HES.8) gives $r\le6$;
the pointed elliptic-curve moduli dimension is at most six (one when $r=0$).
The finite-cover period factorization then gives image dimension below nine.
Thus a putatively nine-dimensional image requires $h=0$, $r\ge12$ and
\[
                  \sum_j\operatorname{rank}(g_j-I)=2e+6.
\tag{HES.9}
\]
This is a necessary condition on the actual image, not a sufficiency claim
from a Hurwitz-space dimension.

The generating tuple has product one. Since $\mu$ is nontrivial, at least
two entries have multiplier different from one. For each such entry the
$1$-eigenspace is totally isotropic for the invariant bilinear form, so its
dimension is at most $e/2$. Those two entries cost at least $e$ in (HES.9).
The other at least ten entries cost at least $e/9$ each. Consequently
\[
                  2e+6\ge e+10e/9,
                  \qquad e\le54.
\tag{HES.10}
\]
If some factor is $S_n$ with $n\ge5$, (HES.2) already gives $e\ge155$,
a contradiction. If all factors are $S_4$, use $e/5$ instead, obtaining
$2e+6\ge e+10e/5=3e$, or $e\le6$, contrary to $e\ge35$.
\end{proof}

\begin{corollary}[D-HES-F: unconditional finite endpoints]
Without assuming (GM), the conclusion of D-HES-03 holds when every simple
factor is one of $\operatorname{Sp}_8(2)$, $\operatorname{Sp}_{10}(2)$ or
$\operatorname{Sp}_{12}(2)$. The same arbitrary scalar, factor-permutation,
and larger-inducing-field scope applies.
\end{corollary}
\begin{proof}
The complete finite tables give projective support at least one fifth of
the degree on every nonidentity core operator. The normalizer and tensor
argument in D-HES-03 therefore gives $e/5$ on every nonscalar inertia.
Every core character is real by the finite input itself. The same unique
bilinear-form and imaginary-character-field argument gives two multiplier
inertias of cost at least $e/2$. The actual intermediate-cover argument
forces $r\ge12$ over genus zero. Hence $2e+6\ge e+10e/5=3e$, or $e\le6$,
contradicting $e\ge35$. Every ingredient in this paragraph is independent
of (GM) and of the higher-rank prime-order argument.
\end{proof}

\paragraph{Source boundary and rejected estimates.}
No small-fourth-moment hypothesis has been added to the inducing
representation. The unrestricted $3/4$ character-ratio bound discussed as
plausible by Gluck--Magaard (2006) has not been used. The proposed stronger
$1/2$ bound is actually false: complete finite tables exhibit ratios $3/5$,
$17/31$ and $11/21$ for $\operatorname{Sp}_8(2)$,
$\operatorname{Sp}_{10}(2)$ and $\operatorname{Sp}_{12}(2)$ respectively.
An initial exploratory GAP comparison using its total ordering on
cyclotomics was rejected because that is not real ordering; the corrected
probe and every retained certificate use integer or rational comparisons
with rigorous absolute-value bounds.
